\subsection{Chilean school admission reports}\label{app:sae}
The Ministry's A1 files give program vacancies; B1 gives the applicant and grade; C1 records the ordered applications with course-level lottery and priority; D1 records the initial offer, before applicants respond. Program means a course within a school, since specialties and schedules can differ at the same school. The 2019 fit uses 483,070 applicants. Application years 2020 and 2021 contribute 454,415 and 461,223 applicants, including those who received no offer. Individuals can apply in successive years; no individual identifier enters the model. Outside the reported list, preferences are unknown. The exercise randomly or systematically prunes a submitted list. It makes no assumption that families supplied truthful rankings over schools they left off that list.

Both fitted rules are ridge logistic edge models trained on whether the initial 2019 offer equals that program. The rank-blind model uses log capacity and demand, a squared congestion term, coarse priority indicators, list length, grade and a smoothed 2019 program offer rate. The last feature is leave-one-out on the training cohort and fixed on test cohorts. Known priorities include returning, sibling, staff-child and former-student statuses. A protected entitlement remains available at every quota and options ranked below it are dropped. The rank-informed model adds log position, an indicator for the first choice and an interaction between position and congestion. Fitted coefficients, exact features and penalty are archived. We retain the highest-scoring $\min(k,d)$ discretionary options and preserve their submitted order. Uniform draws assign an independent random order to the discretionary options. Priorities and lists used for counterfactual allocation are identical across all arms except for policy pruning.

For allocation the benchmark pools vacancies across seat types, ranks applicants by the published priority tier and then by the realized lottery order at that course. At programs with fewer reported vacancies than secured current-school claims, the pooled capacity is raised just enough to accommodate them (6,582 and 6,815 seats across the two years). Grade cohorts are disjoint under fixed priorities. The actual system uses special reserves and family contingencies \citep{correaschool,saereserve}. Its exact published course or outside option agrees with full-list benchmark DA for 432,783 and 440,582 applicants. Official offers below the retained entitlement or off the retained list count as uncovered (3 and 7). C1 appends a current school even when the family did not request it. Excluding placements at such administrative additions from the placement numerator preserves the directions of both five-choice model comparisons. This is a cross-domain test among reported alternatives. Its predictors, population, capacities and use of preference reports differ from those of the rental study.

\begin{table}[t]\centering\small
\caption{Chilean regular-round school admissions. Coverage measures how often an applicant's reference course appears on the five-choice menu plus protected entitlement among applicants assigned by that reference. Placement is the share of all applicants assigned by pooled-seat DA. Uniform reports the mean of 48 independent menu draws per cohort.}
\label{tab:sae}\begin{tabular}{@{}llrrr@{}}\toprule
Cycle & Menu & Official offer (\%) & Full-list DA offer (\%) & DA placed (\%)\\\midrule
2020 & Uniform & 97.59 & 97.57 & 90.66\\
 & First reported choices & 98.77 & 98.83 & 90.34\\
 & Rank-blind prediction & 99.39 & 99.38 & 90.94\\
 & Rank-informed prediction & 99.71 & 99.73 & 90.82\\\midrule
2021 & Uniform & 97.94 & 97.93 & 91.53\\
 & First reported choices & 99.03 & 99.07 & 91.28\\
 & Rank-blind prediction & 99.46 & 99.44 & 91.77\\
 & Rank-informed prediction & 99.79 & 99.80 & 91.66\\\bottomrule\end{tabular}\end{table}

Only 37,312 and 33,570 applicants have more than five discretionary choices. Their rank-blind placement gains relative to uniform were 4.84 and 4.28 points. The other 417,103 and 427,653 applicants experience placement losses of 0.13 and 0.07 points as competitors' lists change. Both official coverage and overall placement increase for the rank-blind model relative to uniform in all 14 grade cohorts in each year. Standard errors over random menus quantify numerical integration only; the 95\% Monte Carlo half-width on the two aggregate placement contrasts is below 0.003 percentage points. Finite-population counts require no independence assumption across applicants. The second ranker and quotas $k=1,2,3$ were specified before scoring, with the comparison between the two fitted models interpreted after observing the primary result. The archive also gives all grade and designated-priority subgroup counts.

As a follow-up, we held the two learned scores fixed and varied the capacity specification. With unadjusted published vacancies, the five-choice rank-blind placement gains relative to uniform were 0.283 and 0.252 points; relative to rank-blind, rank-informed placement fell 0.117 points in both years. With only ordinary vacancies and the same current-school capacity floor, the corresponding paired differences were $+0.281,-0.122$ and $+0.228,-0.114$ points. Neither sensitivity fully implements the deployed reserve routing. A separately implemented simultaneous-offers DA verified every applicant's final program under full exposure and each deterministic five-choice rule, a total of 3,662,552 person--policy checks. For uniform coverage we also summed the exact inclusion probabilities of each reference edge. The Monte Carlo estimates lay within two Monte Carlo standard errors across the sixteen year, reference and quota comparisons.
